\section{Introduction}
\label{sec:intro}

Voice agents are Artificial Intelligence (AI) systems designed to carry out tasks through spoken conversations, and their deployment across a wide range of applications is rapidly growing \cite{moore2025voiceagents}. Voice agents operate under constraints that are fundamentally distinct from text: speech is ephemeral and linear, real-time timing shapes the naturalness of interaction, and acoustic conditions vary widely across callers. These properties give rise to failure modes with no direct text analog \cite{ray2026tauvoicebenchmarkingfullduplexvoice,chen2024voicebench}, and render evaluation frameworks designed for text-based agents \cite{yao2024tau, toolbench2023, patil2025bfcl} insufficient for assessing voice agent quality. Rigorous evaluation of voice agents must therefore address two distinct challenges: how conversations are simulated, and how quality is measured.


The \textbf{simulation challenges} concern constructing interactions that are valid proxies for real deployment conditions. This requires complete multi-turn interactions rather than isolated exchanges - only full conversations expose how an agent recovers from misunderstandings, maintains context across turns, and resolves tasks end-to-end. Conversations must reflect the task-oriented nature of real voice agent deployments, user behavior must reflect natural human spoken dialogue, and acoustic conditions must reflect real-world environments, including variation in accents and background noise. Critically, simulated users must themselves be validated: a simulator that drifts from its assigned scenario, abandons realistic conversational behavior, or acts in ways no plausible human caller would, undermines the validity of any downstream evaluation. Finally, user simulators must behave consistently across repeated runs such that evaluation scores reflect agent behavior rather than simulator variance.


The \textbf{measurement challenge} concerns capturing the full scope of voice agent quality once valid simulations are in place. Task completion and turn-taking dynamics, while necessary, leave critical failure modes undetected \cite{cao2026beyond, mehta2025beyond, andres2025testing}. On the accuracy side, an agent may call the correct tools yet violate system policy, comply with adversarial user requests, or produce spoken outputs containing incorrect entities (e.g. wrong confirmation codes, or monetary amounts) that are catastrophic in production yet undetectable from transcript-level evaluation alone. On the user experience side, an agent may achieve low response latency yet fail to make meaningful progress across turns, repeat prior questions, or present users with an excessive number of spoken options that would overwhelm a user's working memory. Addressing the measurement challenge requires evaluation across a broader set of dimensions than existing benchmarks provide. Additionally, voice agents are not architecturally uniform: cascade systems chain separate speech-to-text (STT), large language model (LLM), and text-to-speech (TTS) components, while audio-native systems process audio inputs directly — either end-to-end via speech-to-speech (S2S) models, or via hybrid systems that pair a large audio language model (LALM) with a TTS model (full definitions in Appendix~\ref{app:definitions}). These architectures have fundamentally different mechanisms, yet must be evaluated on equal footing for benchmarks to meaningfully compare them.


\textbf{We present \framework, a benchmark designed to solve both of these challenges}. On the simulation side, \framework~conducts fully automated bot-to-bot audio simulation over dynamic multi-turn dialogues, with validation-gated quality control ensuring consistency across repeated trials. It includes three enterprise domains comprising 213 scenarios and a perturbation suite of controlled acoustic challenges to probe robustness beyond clean-condition baselines. On the measurement side, \framework~introduces two composite scores: EVA-A (Accuracy) and EVA-X (Experience). EVA-A captures task completion, faithfulness to policy and tool outputs, and audio-level entity fidelity. EVA-X captures conversation progression, conciseness for spoken delivery, and turn-taking timing. Both scores are designed to apply directly to cascade and audio-native architectures, enabling direct comparison across system types. Across 12 evaluated systems, \framework~ reveals that accuracy and experience remain jointly unsatisfied across all architectures, that peak capability consistently overstates reliable performance, and that robustness to acoustic perturbations varies substantially — and non-uniformly — across systems and metrics.
Our contributions are listed below:
\begin{itemize}
\item We introduce \framework: an end-to-end evaluation framework for voice agents that generates realistic bot-to-bot conversations through a user simulator with validation-gated quality control, and supports controlled acoustic perturbations across independent trials.
\item We define EVA-A and EVA-X, joint accuracy and experience metrics that surface failure modes invisible to existing benchmarks and enable direct comparison between audio-native and cascade voice agents.
\item We create three enterprise benchmark datasets with a total of 213 scenarios focused on surfacing voice-specific failure modes.
\item We show empirical findings on cascade vs. audio-native tradeoffs, perturbation sensitivity, and behavioral consistency across trials.
\end{itemize}
